\scanpage{10}
\mspage{3}
% ll. 1–9
Ce morphisme est $W_0$-asphérique, i.e. pour tout $z\in\operatorname{Ob}Z$, son image inverse $f^*(z)=\deleted{\unsure{\Phi|\varphi}}\varphi(z)$ $(=\alpha$, si $z=(\widetilde\Delta_n,\alpha))$ est 0-connexe, ce qui implique que pour tout $G\in\operatorname{Ob}Z^\wedge$ qui est 0-connexe, \deleted{$\unsure{\Phi|\varphi}$} $f^*(\deleted{\unsure{z|F}}G)$ est 0-connexe. (Pour le voir, on \inserted{\unsure{doit|devra|peut}} utiliser le fait que $f^*=\Phi$ commute aux produits fibrés, ou tout au moins aux intersections de sous-objets).
% ll. 9–15
Si on fait appel maintenant à l'hyp. que les objets 0-conn. de $X^\wedge$ (donc aussi ceux des topos induits \unsure{$-/^\wedge$|$Y^\wedge$}) sont 1-connexes, il s'ensuit que $f$ est \deleted{\unsure{un|=}} $\unsure{\mathcal W_1|W_1}$-asphérique, \deleted{i.e.} \inserted{plus précisément} que pour \unsure{tt|$\forall$} \deleted{$z\in\operatorname{Ob}Z$} $G\in\operatorname{Ob}Z^\wedge$ qui est 0-connexe, $f^*(G)$ est 1-connexe.
% ll. 16–21
Mais alors, ayant un morphisme 0-connexe $\deleted{\unsure{X|Y}}Y^\wedge_{/f^*(G)}\longrightarrow Z^\wedge_{/G}$, avec la source 1-connexe, on en conclut que le but $Z^\wedge_{/G}$ l'est aussi. Donc tout objet 0-connexe de $Z^\wedge$ est 1-connexe.
% archive sheet number: 9
